\documentclass[prep,both]{debate}
\motion{对宠物的爱，爱的更是它 / 爱的更是我}
\meta{赛制}{中国科学技术大学新生辩论赛（默认赛制，四对四）}
\meta{生成日期}{2026-09-18}
\begin{document}
\debatetitle
\begin{summary}
\begin{fields}
\claim[持方优劣评估]{\textbf{反方（爱的更是我）较劣势}，劣势在直觉与判准：正方一辩稿把判准钉在「冲突时优先谁」，这把尺子正是反方上一版最想用的，如今正方先拿在手里，还配上了负债、削口粮、照护负担三组行为与损耗数据。反方的补偿在两处：同一份 Lovet 报告里 46\% 因成本延迟就医是行为题，正方引的 91\% 是意愿题；正方引的照护负担研究，同一作者 2022 年证明负担越高越考虑安乐死。反方的打法是全程只做「意愿 vs 行为」「它的指标 vs 我的状况」两个动作，不做道德指控，价值层收在「承认之后才能真的看它一眼」。正方最容易翻的车是把「回报归零」说成绝对，被追问「被需要也是回报」时答不上来；口径是「让人抑郁的被需要，没有人会为自己去追求」。}
\thesis{正方}{钱和它的病相撞时多数人削的是自己；它病了老了不再回应时投入不降反升；明知结局仍然开始。冲突时让路的是我，所以爱的更是它。 }
\thesis{反方}{问卷上愿意负债的是九成，账单真来时让它等的近半；病老之时它一句话说不了，多疼一段还是早走一步，跟着我的负担和不舍走。冲突时让路的是它，决定的是我，所以爱的更是我。 }
\end{fields}
\end{summary}
\section{一、赛制要点}
\begin{flowtable}
\block{A}
\flow{1 正一立论}{180 秒}{正一}{前 20 秒说清「更」是比较、判准是冲突检验；三组数字只报口语形式}
\flow{2 反四质询正一}{120 秒双边}{反四问 / 正一答}{反方要拿到「九成是意愿题」「被需要也是我的需要」两个承认}
\flow{3 反一立论}{180 秒}{反一}{开头一句钉住质询成果；必须先说「我方不否认你爱它」}
\flow{4 正四质询反一}{120 秒双边}{正四问 / 反一答}{正方要逼反方说出「什么行为才算为它」，防判准版定义杀}
\block{B}
\flow{6 反二驳论}{120 秒}{反二}{打正方最难圆的点：负担量表测的是我，同一作者说负担高就考虑安乐死}
\flow{7 正二驳论}{120 秒}{正二}{打反方「意愿 vs 行为」的双标（37\% 也是行为题）；留约 100 字临场位}
\flow{8 二辩对辩}{各 90 秒}{正二 / 反二}{不丢判准，把交锋钉在「有得选时谁让路」}
\block{C}
\flow{10–11 三辩盘问}{各 90 秒}{正三 / 反三}{用「先问二辩再问四辩」检验对方口径}
\flow{12–13 盘问小结}{各 90 秒}{正三 / 反三}{正三在前要预压反三；反三在后留约 80 字临场位}
\block{D}
\flow{15 自由辩论}{各 240 秒}{全员}{三个主战场（钱、病老、终点）推进到位，必问问完，必答不躲}
\block{E}
\flow{17 反四结辩}{210 秒}{反四}{封路式：提前点破正方「回报归零仍在喂药」的画面}
\flow{18 正四结辩}{210 秒}{正四}{留 30–40 秒真回应反四，不背稿}
\end{flowtable}
\textbf{计时方式对策略的影响}：质询是双边计时，这道题的关键承认（「九成是意愿题」「让路发生过」）本身不致命，一辩不要用躲闪换时间，评委看得出来。对辩是单边计时，比的是密度，这道题的对辩会集中在「同一份报告，你引哪一题」，双方二辩要把这一句练到十五秒能说完。自由辩容易滑进「举例子大赛」，谁先把战场钉在「有得选的时刻」谁占便宜。

\section{二、辩题性质与争议焦点}
\begin{fields}
\claim[辩题性质]{\textbf{比较题}（“A 比 B 更……”的变体）。「更」意味着双方都必须承认另一方描述的现象真实存在：正方不能说“这份爱里没有半点自我满足”，反方不能说“人根本不爱动物”。争的是\textbf{同一份爱在两个指向上的权重}。}
\field{正方倾向把题目解释成}{这是一道关于\textbf{让路方向}的题。爱指向谁，看它的需要与我的需要相撞时谁让路；而最能读出方向的，是它什么都给不了我的那一段。正方希望评委在“冲突时刻你做了什么”这个层面上判断。}
\field{反方倾向解释成}{这是一道关于\textbf{决定变量}的题。同一个冲突时刻，让它等还是让它走，跟着的是我的钱、我的负担、我的不舍；它的指标只是被我读出来的。反方希望评委在“行为数据说什么、决定跟着谁走”这个层面上判断。}
\field{争议焦点}{三个战场双方必然相遇：（1）\textbf{钱}：同一份 2026 年两千人报告，91\% 愿负债与 37\% 削口粮是正方的，46\% 延迟就医是反方的，争的是意愿题还是行为题、付不起算不算让路；（2）\textbf{病老}：238 人对照研究里照护者负担与抑郁显著更高，正方读成“回报归零投入最高”，反方读成“负担长在我身上，负担高就考虑安乐死”；（3）\textbf{终点}：安乐死的理由是它的疼痛（48.5\%、87\%），还是我的负担与依恋（中介效应显著、依恋是主导因素 69.3\%）。}
\end{fields}

\section{三、关键词定义}
\subsection{3.1 爱}
\begin{fields}
\field{调研}{\sub{词典义}{「爱」第一义项“对人或事有深挚的感情”（《现代汉语词典》体系，汉典 zdic.net，2026-09-15 访问）。汉语的“爱”同时覆盖感情状态与关切行为（爱护、爱惜），这个双重性是本题的定义战场。}\sub{学术义}{Stanford Encyclopedia of Philosophy「Love」条（2026-09-18 访问）。稳健关切说：Frankfurt（1999, p.129）说爱“neither affective nor cognitive. It is volitional”，是一套为被爱者自身之故的稳定动机结构。结合说：Nozick（1989）说恋人把各自的福祉“pooling”成一个共同池。非替代性：Badhwar（2003, p.63）说被爱者“phenomenologically non-fungible”。历史关系属性：Whiting（1991）、Delaney（1996）用共同经历解释为什么不换一个条件相同的人。}}
\begin{procard}{正方有利定义}\definition{\textbf{爱是为对象自身之故而生的关切与行动}（稳健关切说）。}\sub{有利之处}{一旦把“为对方之故”写进定义，凡是真正的爱就必然指向对方，反方要么承认多数人对宠物的感情不是爱，要么承认正方立场。}\sub{反方如何反驳它}{标准的\textbf{定义杀}，把结论藏进了定义。按这个定义“爱的更是我”在语义上不可能成立，辩论失去意义。反方当场反问：“按对方的定义，一个知道狗超重却继续喂零食的主人，算不算爱狗？”}\end{procard}
\begin{concard}{反方有利定义}\definition{\textbf{爱是一种满足主体心理需要的情感投入}。}\sub{有利之处}{爱被定位成“我的一种状态”，归属天然在我这边，一切为动物做的事都可以重述为“我在满足我的情感需要”。}\sub{正方如何反驳它}{同样是定义杀，而且让命题\textbf{不可证伪}：母亲对孩子的爱、战友之间的爱全都“更是我”。正方直接问：“按您的定义，有没有任何一种爱是更指向对方的？”反方必须答“有”，然后要条件。}\end{concard}
\begin{neutralcard}{中立定义}\textbf{爱是对某个对象产生的深挚感情，以及由这份感情带来的持续关切与行动。}它有来源（因为什么而起）、有内容（想要对方怎么样）、有代价（愿意为它承担什么）。\sub{合理性}{贴近词典第一义项，把“感情”与“关切行为”两层都保留；不预设指向，双方都有得打。这一版双方一辩稿都用，不争。}\end{neutralcard}
\end{fields}
\subsection{3.2 更是}
\begin{fields}
\field{调研}{「更」在比较句中表示程度上的进一步，预设比较项同时存在。“A 更是 X”的标准读法是“在 X 与 Y 之间 X 占的分量更重”，不是“只有 X”。两方一辩稿都要说清这一点，否则会被对方指为绝对化。}
\begin{procard}{正方有利定义}\definition{\textbf{「更」看冲突时刻的让路方向，尤其是回报归零的时刻投入是否还在。}}\sub{有利之处}{把比较收窄到最能读出方向的场景，日常里同向的部分不计分；正方的负债、削口粮、照护负担数据全部落在这里。}\sub{反方如何反驳它}{只看最感人的那一段，等于用 5\% 的时刻判整段关系；而且“让路”本身可能是为了避免自己失去之痛，让路不等于为它。}\end{procard}
\begin{concard}{反方有利定义}\definition{\textbf{「更」看这份爱的受益结构与内容由谁塑造。}}\sub{有利之处}{挑品种、喂法、拍照、当孩子，内容全由我定；受益结构里陪伴、被需要、自我形象全在我。}\sub{正方如何反驳它}{受益结构标准下母爱、爱情全部“更是我”，不可证伪；内容塑造回答的是“怎么爱”，不是“更爱谁”。}\end{concard}
\begin{neutralcard}{中立定义}\textbf{「更」，就是当它的需要和我的需要分离、相撞的时候，这份爱让谁受益、让谁让路；以多数饲主的典型表现为准，极端个例只用来说明机制。}\sub{合理性}{比较题需要一把能分出高下的尺子，日常里同向的部分分不出双方，冲突处才看得见指向；两方都有硬数据可打（正方：负债、削口粮、照护负担；反方：延迟就医、被告知后不改、负担预测安乐死考虑）。正方一辩稿把它说成“冲突时优先谁”，反方一辩稿用同一把尺子。}\end{neutralcard}
\end{fields}
\subsection{3.3 它 / 我}
\begin{fields}
\field{它}{\textbf{你家那一只。}它有自己的生理与行为需要，不随你的想法改变；它会区分人，是一个真实的他者。}
\field{我}{\textbf{饲主自己的需要}：陪伴、被需要、情感出口、可控感、自我形象。反方全场不把它说成“自私”，正方也不能把反方的“我”说成道德指控。}
\field{时间范围}{“对宠物的爱”从它进家门、成为“这一只”开始算，到它离世为止。挑品种是买之前的事（正方口径），反方则主张明知会喘还买、进家之后否认它在喘，选择每天都在重复。离世之后的葬礼、克隆不在冲突检验的范围里，只能作画面，不能作证据。}
\end{fields}
\subsection{3.4 隐含要素}
\begin{fields}
\field{主体}{中国当下普通城镇饲主（1.26 亿只犬猫背后的家庭），不含商业繁育者。双方的多数比例数据大多来自英美，中国数据用来证明账单真实与机制存在，这一点双方对等，谁先承认谁占便宜。}\field{范围}{日常养宠关系里的三个冲突时刻：钱、病老、终点。}\field{时间尺度}{从进家到离世的整段关系，重点是病、老、死三个回报归零的阶段。}\field{比较对象}{同一份爱在两个指向上的权重，不是“爱它”与“不爱它”。}
\end{fields}
\begin{warnbox}{3.5 定义杀陷阱与判准版定义杀：双方都要背}
\begin{fields}
\begin{procard}{正方的陷阱}不要把“爱”定义成“为对方之故”；更不要把每一种行为都读成“为它”（付钱是为它、停钱也是为它、撑着也是为它）。反方一定会问：“在你方立论下，饲主做什么才算为我？”正方必须当场给出三样：有得选时让它扛、拖着不让它走、被告知之后仍不改。然后说：这三样都存在，我方证的是它们是少数。\end{procard}
\begin{concard}{反方的陷阱}不要把“爱”定义成“满足自我需要”；更不要说“安乐死是它让路、继续治也是它让路”。正方一定会问：“在你方立论下，饲主做什么才算为它？”反方必须当场给出三样：有得选时让自己扛、按它的指标而不是我的状况定终点、被告知之后按它的需要改。然后说：这三样都存在，我方证的是多数人跟着自己的状况走。\end{concard}
\end{fields}
\end{warnbox}

\section{四、判准与论证义务}
\begin{fields}
\claim[判准是什么]{评委凭\textbf{冲突检验}判断：当它的需要与我的需要分离甚至相撞时，这份爱让谁受益、让谁让路。这是正方一辩稿明说的“冲突时优先谁”，反方一辩稿也接受这把尺子。}
\begin{procard}{正方有利判准}\definition{只看冲突时刻的让路方向，尤其是回报归零时投入是否还在。}\sub{有利之处}{正方三组数据（负债与削口粮、照护负担、终点按它的疼痛）全部落在这里；“回报归零投入最高”一句话就能把“图回报”解释排除。}\sub{反方如何反驳它}{“回报归零”在共同定义下不成立：被需要、可控感、自我形象在病老期恰恰最高；把日常全部排除等于只用最感人的时刻判整段关系；让路可能只是避免我的失去之痛。}\end{procard}
\begin{concard}{反方有利判准}\definition{看这份爱的受益结构与内容由谁塑造。}\sub{有利之处}{陪伴、被需要、自我形象全在我；挑品种、喂法、拍照全由我定。}\sub{正方如何反驳它}{受益结构标准让母爱、爱情全都“更是我”，不可证伪；内容塑造回答的是“怎么爱”不是“更爱谁”。请反方回到冲突检验。}\end{concard}
\begin{neutralcard}{中立判准}\textbf{冲突检验：当它的需要和我的需要分离或相撞时，这份爱让谁受益、让谁让路；以多数饲主的典型表现为准，极端个例只用来说明机制。}\sub{合理性}{符合比较题性质（要一把能分高下的尺子）；照顾题目主体（普通饲主的三个必经时刻）；能真正区分双方（正方要证多数人让自己，反方要证多数人让它或决定跟着我走）。}\sub{正方需要证明}{在真实冲突里，典型饲主让的是自己；并且在回报最低的阶段投入不降反升，“图回报”解释不了这份投入。}\sub{反方需要证明}{在真实冲突里，典型饲主让的是它；或者所谓“让路”跟着的是我的状况（钱、负担、不舍），而不是它的指标。}\end{neutralcard}
\field{一辩稿用哪一版}{双方一辩稿都用中立判准。正方的一辩稿已经把它说成“冲突时优先谁”，反方一辩稿接受并补一句“以行为为准、以多数为准”。有利判准留在攻防里用：对方偏向时指出偏向，再把评委拉回中立判准。}
\end{fields}
\prosection{五、正方论点（爱的更是它）}
正方一辩稿已由队伍定稿，三个论点按稿件顺序编号。全队三句口径：（1）判准是冲突时优先谁，说全两问：让谁受益、让谁让路；（2）它的疼通过我的负担到达我的决定，渠道是我，起点和理由是它；（3）被需要是我的需要，它需要我起床、我需要睡觉，起床的是我。
\prosubsection{论点一：回报归零，投入最高}
\begin{fields}
\claim{它病了、老了、不再回应，它能给我的陪伴、互动、回应归零；典型饲主的投入不降反升，并在自己身上留下可测量的伤害。图回报，解释不了这件事。}
\begin{mechanism}\step 慢性病、终末期、瘫痪让它能给我的趋近于零，而它的需要每天和我的睡眠、收入、健康相撞：凌晨三点它叫，我起不起。\step 这一刻退路最便宜：送救助站、不治、提前安乐死都在手边；可控感、自我形象都留不住人，留下来只能是为它。\step 典型饲主继续照护，直到自己出现临床水平的负担与抑郁；典型的狗在诊所里被养到中位十二岁多。\end{mechanism}
\begin{evidence}{学理}[已核实]稳健关切说（robust concern view），Harry Frankfurt，1999，《Necessity, Volition, and Love》p.129；2004《The Reasons of Love》系统化。\sub{核心主张}{爱“既非情感也非认知，而是意志性的”（neither affective nor cognitive. It is volitional），是为所爱者自身之故的稳定关切，其持续不以对方回报为条件。}\sub{支撑}{第 3 步：解释为什么回报归零投入不降，动力在承诺不在回报回路。}\sub{边界}{只用来解释数据，不用来裁定辩题（否则是定义杀）。反方会说“守承诺是为了自我形象”，答：判准问让谁让路，守承诺让出去的是睡眠、收入、健康。}\sub{出处}{Stanford Encyclopedia of Philosophy「Love」条。}\end{evidence}
\begin{evidence}{数据}[已核实]238 位猫狗饲主对照研究：119 位照护慢性或终末期疾病动物 vs 119 位匹配对照；照护组负担、压力、抑郁与焦虑症状更高，生活质量更差，全部 p<0.001。\sub{来源}{Spitznagel MB 等，2017，Veterinary Record 181(12):321，DOI 10.1136/vr.104295；作者在 Kent State 新闻稿里说“在许多方面与照护患痴呆家人的情况相似”。}\sub{方法}{横截面对照，标准化量表（Zarit 负担量表等），美国网络招募，按饲主年龄性别与动物物种匹配。}\sub{局限}{自报、横截面，样本是“正在照护”的人。它证明的是照护者仍在照护并且在受损；“投入”的定义就是让出去的睡眠、时间、钱、健康，负担量表测的正是这个。反方会拿同一作者 2022 年“负担高就考虑安乐死”反打，口径见第 2 条数据。}\end{evidence}
\begin{evidence}{数据}[已核实]同一作者 2021 年 393 名年老或重病伴侣动物饲主：照护负担、预期性哀伤、生活质量三者不可互换；2022 年 277 名骨关节炎犬主人：它的疼痛与功能受损 → 我的负担 → 考虑安乐死，中介效应显著（P≤0.01）。\sub{来源}{Spitznagel 等 2021，Vet Rec 188(9):e74；Spitznagel 等 2022，Vet J 286:105868。}\sub{口径}{反方一定会念 2022 年这条。正方的回答只有一句：它的疼通过我的负担到达我的决定，渠道是我，起点和理由是它。负担影响决定有多难，不改变决定的内容是它的痛（理由 87\%、48.5\%）。绝不能说“我的负担不决定它哪天走”，那句和这两篇论文原句冲突。}\sub{局限}{横截面；“考虑”不是“做”，277 人都还在照护是抽样条件不是发现。}\end{evidence}
\begin{evidence}{数据}[已核实]67 名瘫痪或轻瘫犬主人 vs 90 名健康犬主人：负担显著更高（p<0.001）、生活质量更低；英国 VetCompass 近三万例确认死亡中 89.3\% 为安乐死、8.3\% 自然死亡，安乐死犬中位年龄 12.1 岁；美国两千人调查中 71\% 说宠物的未来一直在自己的担忧里。\sub{来源}{Pryjmak 等 2026，Front Vet Sci；Pegram 等 2021，Sci Rep 11:9145；Talker Research 为 MetLife 2026-01-30。}\sub{用法}{“多数”的落点在这里：典型的狗被养到老、养到病、养到走，不是在病老那一刻被送走。VetCompass 只量记录在案的死亡，不用它说弃养比例。}\end{evidence}
\begin{evidence}{案例}[已核实]慧慧与哈士奇卢卡：15 年前带回家；13 岁起失禁、器官衰竭、瘫痪、痴呆；“最高纪录是一晚上起床六七次，第二天再去正常上班”；2024 年辞职专心照顾，每 3 小时回家一次；14.9 岁“因无法自主吞咽，不得不接受安乐死”。\sub{代表性}{主流媒体深度报道、中国青年饲主。典型在数据里（p<0.001 的整组照护者），案例只给画面。反方会说“上新闻所以不典型”“两年失禁瘫痪痴呆就是拖着”：答，吞咽是它的指标；反方自己承认按它的指标放手是为它。不要加原文没有的细节（例如“和兽医一起判断”）。}\sub{出处}{新京报 秦冰 2025-05-24《老年宠物，年轻人的爱与“不可承受之重”》。备用：郑州吴先生为瘫痪八年老狗配轮椅，澎湃新闻 2020-05-08。}\end{evidence}
\closing{病老是回报最低、退路最便宜的冲突时刻；受益的是它（被照护），让路的是我的睡眠、收入、健康。}
\begin{clash}
\pair{被需要、可控感、自我形象在病老期恰恰最高，回报没有归零。}{被需要是我的需要，我方承认。判准问相撞时谁让路：它需要我起床，我需要睡觉，起床的是我。被需要得到满足，是因为我让了路。可控感：它失禁、痴呆、随时会走，我控制的只是自己要不要起床。}
\pair{负担量表测的是代价落在我身上，代价不说明方向；同一作者说负担高就考虑安乐死。}{判准两问：让谁受益、让谁让路。代价落在谁身上，谁就让了路；受益的是被照护的它。E16 的路径起点是它的疼，渠道是我的负担，做出的决定理由是它的疼。}
\pair{它的处境跟着我的状况走，我的状况不随它走。}{我的状况也跟着它的处境走，p 小于千分之一（238 人、157 人、277 人三组）。关系是对称的，对称就只能看谁让路，病老期让路的是我。}
\end{clash}
\end{fields}
\prosubsection{论点二：让路的是我}
\begin{fields}
\claim{钱和它的病相撞时，多数人削的是自己：削减自己的口粮去付它的药费。这就是“冲突时优先谁”最朴素的答案。}
\begin{mechanism}\step 它的病变成账单，和我的伙食、房租、休息正面相撞。\step 退路便宜：少治、不治、送走都在手边，代价对我很低、对它很高。\step 九成人说愿意为救它负债；三成七的人已经在过去一年削了自己的食品、水电和账单，三成七已经真的负了债，四成四差点因钱放弃但最后付了。\end{mechanism}
\begin{evidence}{学理}[已核实]共情—利他假说（empathy–altruism hypothesis），C. D. Batson 等，1981，J Pers Soc Psychol 40(2):290–302。\sub{核心主张}{利他动机与利己动机的区别，在于“退路便宜的时候还帮不帮”：高共情组逃离容易时帮助率 .91、困难时 .82，不受退路影响。}\sub{支撑}{第 2、3 步的判别逻辑：养宠里最便宜的退路是不治和送走，多数人不走。}\sub{边界}{人对人实验室研究，每格 11 人，只提供判别逻辑，比例由数据给。反方会说退路不便宜（罚款 500 元、自我形象、失去快乐来源）：答，判准问让谁让路，不问让路时心里舒不舒服；舍不得而削自己的口粮，让路的是我。}\end{evidence}
\begin{evidence}{数据}[已核实]2026 年 2,000 名美国宠物主：91\% 愿为救宠物负债；37\% 过去一年用信贷或负债付兽医费；37\% 削减食品杂货、水电、账单；44\% 上次就医差点因成本放弃；46\% 因成本延迟或跳过就医；16\% 遇 2,000 美元急诊会延迟或放弃。\sub{来源}{Lovet《2026 Checkup Report》，Censuswide，2026-05-21 至 27。}\sub{方法}{在线自报问卷。91\% 与 16\% 是意愿与假设题，37\%、37\%、44\%、46\% 是行为题。}\sub{局限}{宠物医疗融资公司委托；美国样本；同一份报告反方引 46\%。\textbf{诚实口径}：行为题两边都不过半，人数相当。正方不在钱上争多数，只争方向；压倒性的证据在论点一和论点三。不要用 84\%（16\% 的反面）证多数，那是假设题。}\end{evidence}
\begin{evidence}{数据}[已核实]2,000 名美国犬猫主：65\% 会先削减自己的生活方式预算再动宠物的；Rover 1,000 人：33\% 已在其他生活领域减支以保证宠物所需、34\% 说宠物支出是最后砍的类别之一；中国单只犬年均 3006 元、医疗占 27.6\%；狗跟腱手术 5500 元、猫 CT 8000 元。\sub{来源}{Talker Research 为 MetLife 2026-01-30；Rover 2025-03-18；派读白皮书 2026-01-05；HUST 数说 2025-11-28。}\sub{用法}{中国数据只证账单真实，不证方向；65\% 是意愿题只说方向。}\end{evidence}
\begin{evidence}{案例}[已核实]2020 年武汉封城：求援者手持身份证录视频，委托素不相识的志愿者开锁进自己家；8 人协会加 60 名志愿者、约 20 名开锁师傅，1 月 26 日至 2 月 11 日为将近 5000 只留守宠物补充食水。\sub{代表性}{极端情境，只说机制不证比例。口径：主人让出去的是隐私和安全感，换它一口水。反方反用“出门时把它留下”：封城是 1 月 23 日凌晨突然宣布的，主人是回不去。}\sub{出处}{腾讯新闻谷雨工作室 2020-02-28。}\end{evidence}
\closing{钱的战场方向清楚、人数相当：问卷里说的是为它，行为里做的也有让自己；反方拿到的只是“有人让它等过”。}
\begin{clash}
\pair{同一份报告：91\% 是意愿题，46\% 是行为题，行为跟着我的钱走。}{37\%、37\%、44\% 也是行为题，让的是自己；行为题两边都不过半。反方要的“多数让它等”同样没有。}
\pair{付不起就让它等，这就是相撞，就是让路。}{接受。付不起的人一部分借钱削口粮，一部分让它等，谁也没证明多数。}
\pair{Batson 的退路在养宠里不便宜，不走退路是自利模型的预测。}{判准问让谁让路，不问让路时舒不舒服。舍不得而削自己的口粮，让路的还是我。}
\end{clash}
\end{fields}
\prosubsection{论点三：明知结局，仍然开始（终点按它来定）}
\begin{fields}
\claim{饲主明知它会先走仍然开始；到终点那一步，多数人按它的痛苦、不按自己的不舍做决定，主动承受自己一生里最痛的失去之一。}
\begin{mechanism}\step 它是不可替代的“这一只”，失去它是真实的失去：21\% 把宠物之死列为一生最痛，7.5\% 达到延长哀伤障碍标准。\step 终点上“多留它一天”与“让它少疼一天”正面相撞，这是全场最纯粹的冲突。\step 多数人按它来定：安乐死理由是疼痛与痛苦（48.5\%，646 人；87\%，8,955 人），疾病与衰退（63.3\% 加 31.6\%）；依据是观察它（53.5\%）和听兽医（39.2\%）。\end{mechanism}
\begin{evidence}{学理}[已核实]非替代性论（phenomenological non-fungibility），Neera Badhwar，2003，p.63，SEP「Love」条转引。\sub{核心主张}{被爱者在现象上不可替换。}\sub{支撑}{第 1 步：如果爱更是我，换一只同款即可续上我的需要，剧痛不应发生；剧痛发生且被预见，说明对象不可换。}\sub{边界}{非替代性双方都承认，不单独判方向；只用它解释为什么痛。}\end{evidence}
\begin{evidence}{数据}[已核实]英国 975 人全国配额样本：32.6\% 经历宠物死亡，其中 21.0\% 列为一生最痛的失去；延长哀伤障碍条件发生率 7.5\%。中国青年饲主 71.6\% 困惑“它老了怎么办”。\sub{来源}{Hyland 2026，PLOS ONE；青年志 2025-02-28。}\sub{用法}{71.6\% 只用来说“开始时就知道结局”，不说担忧的内容。}\end{evidence}
\begin{evidence}{数据}[已核实]Dog Aging Project 终末调查 646 名犬主：83.0\% 安乐死，首要理由疼痛或痛苦 48.5\%、生活质量差 24.8\%、预后差 19.6\%；同一项目 8,955 名犬主中 87\% 把疼痛或痛苦列为理由，近期兽医就诊与感知疼痛显著相关；228 名宠物主（17 国）：理由疾病 63.3\%、衰退 31.6\%，依据自己观察 53.5\%、兽医建议 39.2\%、结构化生活质量工具 7.3\%，情感依恋是决策主导因素 69.3\%。\sub{来源}{McNulty 等 2026，JAVMA 264(7)；Mann 等 2026，AJVR；Kiss 等 2026，Animals 16(11):1738。}\sub{局限与口径}{全是自报，反方的数据也是。反方会引同一摘要的 69.3\%（依恋主导）和 7.3\%（结构化工具）：答，69.3\% 不带符号，依恋深的人可能拖也可能不忍它再疼而早放手；观察的对象是它，兽医看的也是它；反方要证“感知随依恋变”，没有这个数据。“所有疼痛都是被感知的”这句可以说，但要接上“观察的对象是它”。}\end{evidence}
\begin{evidence}{案例}[已核实]卢卡 14.9 岁不能自主吞咽，慧慧“不得不接受安乐死”。\sub{代表性}{与论点一同一案例的终点，场上一线串起：为它辞职，为它起床，最后也为它放手。吞咽是它的指标。}\sub{出处}{新京报 2025-05-24。}\end{evidence}
\closing{死是回报永久归零、我的痛最大的冲突时刻；多数人按它的痛来定，让路的是我，受益的是它。}
\begin{clash}
\pair{没有人为爱自己选剧痛？人会为自己结婚生子，85\% 说它是主要快乐来源，剧痛是快乐的账单。}{承认，这句在一辩稿里是价值收束，不是证据。证据在终点：签账的人在快乐归零那天可以止损，多数人没有。}
\pair{理由是事后填的，变量是当时起作用的：依恋 69.3\%、我的负担中介。}{同一张问卷，理由与因素同为自报；带方向的数字（63.3\%、31.6\%、39.2\%）全在正方，69.3\% 不带符号。}
\pair{确信时机正确的人仍有 47\% 内疚，内疚是“我是不是好主人”，是自我形象。}{内疚是亏欠了谁的情绪，人不会为亏欠自己而内疚；“好主人”的对象是它。25\% 后悔对应 75\% 不后悔。}
\end{clash}
\end{fields}
\subsection{（被淘汰的正方候选论点，交给反驳库与自由辩备用）}
\begin{enumerate}[leftmargin=1.8em]
\item 极端情境的让路（封城、灾害）：与论点二同机制，并入作案例。
\item 哀伤不被承认仍在哀伤（57.1\% 必须隐藏哀伤）：冲突已结束，判准无从检验；反方一句“痛的是我”即可打掉。只作反驳库回应。
\item 意志性关切（Frankfurt、Badhwar）：纯学理，单独立论是定义杀；拆入论点一、三作学理。
\item 日常作息为它让路（遛狗、居住、出行）：无可核实比例；反方一句“习惯就是我的生活”打掉。
\item 它认得我、它爱我（陌生情境实验）：因变量在狗身上，一说就被反方接走；只在反驳库里用“真实他者”口径。
\end{enumerate}

\consection{六、反方论点（爱的更是我）}
反方在同一条判准下作战，两个论点统一在一句命题下：冲突时让不让路，跟着的是我的钱、我的负担、我的不舍，而不是它的指标。全队三句口径：（1）我方不否认你爱它，我方比的是分量；（2）同一份报告，正方引意愿题，我方引行为题；（3）病老终点它没有发言权，决定它多疼一段还是早走一步的，是我的负担和我的不舍。
\consubsection{论点一：说的是为它，做的跟着我的钱走}
\begin{fields}
\claim{钱和它的病相撞时，它等不等，由我的余额决定；“愿意负债”是我对自己的叙述，“拖过就医”是我做的事。}
\begin{mechanism}\step 同一份两千人报告：意愿题九成一愿为救它负债；行为题四成六过去一年因成本拖过或跳过就医、四成四上次差点因钱放弃、两成五改买更便宜的粮。\step 意愿与行为之间的差距，心理学叫意图—行为鸿沟：意愿的大幅变化只带来行为的小幅变化。\step 判准看相撞那一刻让谁让路：账单来了，先等的是它的病；说的是为它，做的跟着我的钱走。\end{mechanism}
\begin{evidence}{学理}[已核实]意图—行为鸿沟（intention–behavior gap），Sheeran 2002；Sheeran \& Webb 2016，Social and Personality Psychology Compass 10(9):503–518。\sub{核心主张}{操纵意愿的实验元分析显示，意愿的中到大幅变化只带来行为的小到中幅变化（d+ = .36），鸿沟主要来自“有意愿却不行动的人”。}\sub{支撑}{第 2 步：要求评委给行为题更高权重。}\sub{边界}{只说明意愿不等于行为，不指控饲主说谎。不要把 16\%（急诊假设题）和 46\%（全年行为题）相减说成“三十个点的鸿沟”，两题不同。}\end{evidence}
\begin{evidence}{数据}[已核实]2026 年 2,000 名美国宠物主：46\% 因成本延迟或跳过就医；44\% 上次就医差点因成本放弃；16\% 遇 2,000 美元急诊会延迟或放弃；91\% 愿负债；37\% 削减食品水电账单；37\% 负债。Rover 1,000 人：25\% 改买更便宜的宠物食品或服务。\sub{来源}{Lovet《2026 Checkup Report》，Censuswide；Rover 2025。}\sub{方法}{在线自报问卷。}\sub{局限与口径}{企业委托、美国样本；与正方同一份报告，先说“意愿题他们引，行为题我们引”。\textbf{诚实口径}：行为题两边都不过半（削自己三成七、让它等过四成六），反方不说“多数让它”，反方说正方所有过半的数字全是意愿和假设题，那是人怎么说自己，不是人怎么做。不要说“46 大于 37”，两题不互斥。}\end{evidence}
\begin{evidence}{数据}[已核实]中国单只犬年均 3006 元、猫 2085 元，医疗占 27.6\%；狗跟腱断裂手术 5500 元、猫 CT 8000 元。\sub{来源}{派读白皮书 2026-01-05；HUST 数说 2025-11-28。}\sub{用法}{证明中国语境下“账单来了”的量级，冲突不是假想。}\end{evidence}
\begin{evidence}{案例}[已核实]北京房山救助基地八哥犬“月饼”：2024 年中秋当天送来，“主人就联系我们，说自己要去外地，而家里其他人都不想继续养犬了”。\sub{代表性}{只证“生活变动时被移走的是它”这一机制真实存在，不证比例。正方说“多数人没弃养”：反方不比弃养人数，比的是冲突真的发生时方向跟着谁。}\sub{出处}{中国日报网转法治日报 2025-01-23。备用：猫“懒懒”患猫癣和肾积水后被家人扔在救助站附近（青年志 2023-10-13）。}\end{evidence}
\closing{钱的战场，行为题打平；让不让它等由我的余额定；“为它”是叙述，“让它等”是行为。}
\begin{clash}
\pair{37\%、37\%、44\% 也是行为题，让的是自己；行为题两边都不过半。}{承认。所以正方“多数人让自己”没有证据；正方所有过半的数字是意愿和假设题。行为题上评委判平，方向看决定变量。}
\pair{付不起不是相撞，是匮乏；有得选的人在急诊档 84\% 不拖。}{我的房租、伙食也是需要，付不起就是相撞，先等的是它。84\% 是假设题，和九成一同一题型，正方自己说假设题不证多数。}
\pair{“嘴上说为它”本身就是感情的方向。}{共同定义里“我”的第五项是自我形象；问卷上说“我会为它负债”，说的正是我想成为的那种主人。}
\end{clash}
\end{fields}
\consubsection{论点二：病老终点，舍不得就让它多受一段}
\begin{fields}
\claim{正方自己挑的“回报归零”考场上，它没有发言权；决定它治不治、走不走的变量是我的负担和我的不舍，而不舍的结果是它多受一段。}
\begin{mechanism}\step 病老临终时决定权全在我。\step 决定变量是我的：决策主导因素是情感依恋（69.3\%），依据是自己观察（53.5\%），用结构化生活质量工具的只有 7.3\%；同样的疼痛，我的负担决定我考不考虑放手（中介效应显著）。\step 不舍的方向是让它多受：79\% 的兽医被要求过提供自己认为无效的治疗，76.2\% 的兽医承认无效治疗有时有益的是主人；主人事后 53\% 内疚，确信时机正确仍有 47\% 内疚。\end{mechanism}
\begin{evidence}{学理}[有把握]照护负担理论与 Zarit 负担量表，Zarit, Reever \& Bach-Peterson，1980，The Gerontologist 20(6)。\sub{核心主张}{负担是照护者对照护如何影响“自己生活”的主观评价；Spitznagel 2017、2022 把它移植到宠物照护。}\sub{支撑}{第 2 步：决定它生死的变量之一是我的负担分。}\sub{边界}{正方会问“这个量表是为照护失智父母的子女设计的，他们也更爱自己吗”：答，反方不用负担分判母爱，用它判“决定跟着谁走”；失智父母送不送机构子女也跟着自己的负担走，没有人因此说他们不爱父母，反方也没有说饲主不爱它，反方说的是分量。补充学理：Ellis 2024（1,359 人）依恋越高抑郁焦虑孤独越高，病老期投入护住的是我的依恋。}\end{evidence}
\begin{evidence}{数据}[已核实]228 名宠物主（17 国）：情感依恋是决策主导因素（69.3\%）；依据自己观察 53.5\%、兽医建议 39.2\%、结构化生活质量工具 7.3\%。277 名骨关节炎犬主人：照护负担对“考虑安乐死”有显著间接效应。\sub{来源}{Kiss 等 2026，Animals 16(11):1738；Spitznagel 等 2022，Vet J 286:105868。}\sub{局限与口径}{同一份 Kiss 摘要里理由是疾病 63.3\%、衰退 31.6\%，正方会引；反方的口径是“理由是事后的解释，主导因素是当时的变量，两栏都是自报，反方不说哪个假”。正方说“负担是它的疼的函数”：反方不否认我的状况随它的病变，反方说决定跟着的是“它的病落在我身上的量”，不是它的指标本身，用指标的只有 7.3\%。}\end{evidence}
\begin{evidence}{数据}[已核实]889 名北美兽医：79\% 被要求提供自己认为无效的治疗；477 名小动物兽医：99.2\% 职业生涯遇到过、85.0\% 过去一年遇到过；作者评论文：76.2\% 同意无效治疗有时有益的是主人。英国宠物主调查：53\% 宠物死后内疚，确信时机正确仍 47\% 内疚，25\% 后悔。\sub{来源}{Moses 等 2018，JVIM 32(6)；Peterson 等 2022，JAVMA 260(12)；Peterson \& Boyd，Hastings Center 2022-06-21；Vets at Home × Pet Database，Companion Life 2026-09-07。}\sub{局限与口径}{兽医数据是兽医被要求过的比例，不是饲主比例，只证“舍不得就让它多治一段”是普遍现象；76.2\% 出自作者评论文，引用时注明；英国调查样本量未公布，只说“英国一项调查”。}\end{evidence}
\begin{evidence}{案例}[已核实]卢卡 13 岁起失禁、瘫痪、痴呆，14.9 岁不能自主吞咽才“不得不接受安乐死”，中间将近两年。\sub{代表性}{正方自己的案例。用“还能不能活”当开关而不用“过得好不好”（正方引的 Dog Aging Project 里生活质量差是第二大理由），反方不评价对错，只指出开关跟着“我什么时候接受得了”走。结辩画面：德牧“包包”安乐死后克隆，预计花费 30 万以上，它已经不在了，30 万买的是我的不失去（离世之后，不作证据）。}\sub{出处}{新京报 2025-05-24；青年志 2023-10-13。}\end{evidence}
\closing{回报归零的阶段，决定跟着我的依恋和负担走，不舍的结果是它多受一段；让路的是它，决定的是我。}
\begin{clash}
\pair{E16 起点是它的疼；做出的决定理由是它的疼（87\%、48.5\%）；开关在它的指标上。}{同样的疼落在不同的我身上，负担分不同、决定不同；用结构化指标的只有 7.3\%，“它的指标”是我读出来的。}
\pair{89.3\% 安乐死、8.3\% 自然死，拖到死的是少数。}{安乐死之前拖了多久 VetCompass 没有；有比例的是兽医被要求无效治疗的普遍程度与主人的内疚。反方承认没有饲主端拖延时长比例，正方也没有“按指标及时放手”的比例。}
\pair{一说钱紧让它等、二说舍不得花钱多治，典型饲主是两个人。}{同一个人：钱紧时它等，舍不得时它多治，两次跟着的都是我的状况；正方的典型饲主要同时满足“付得起就付、指标一到就放手”，两件事正方都没有行为比例。}
\end{clash}
\end{fields}
\subsection{（被淘汰或降级的反方候选论点，交给反驳库与自由辩备用）}
\begin{enumerate}[leftmargin=1.8em]
\item 知道了也不改（喂胖、穿衣、当孩子）：降级为自由辩追问。61\% “了解但无应对”的分母是全体饲主，中国犬猫超重只有两到三成，证不了多数；短头犬主人 58\% 否认它喘只有 31 例，是案例；穿衣致热休克只有“可致”没有发生率。追问句：“它要少吃，你要它吃得开心；它要散热，你要它穿得可爱。相撞了，让的是谁？”
\item 让路的报酬仍在我（受益结构）：正方一句“母爱也更是我，不可证伪”打掉；证据并入论点二。
\item 生活一变就弃养：无可靠比例，“流浪动物 4000 万”未核实；个案并入论点一。
\item 审美选择每天重复（挑品种）：触发时间范围争议；58\% 并入追问。
\item 它老了怎么办是担忧我的未来（71.6\%）：正方一句“担忧它的老年正说明爱它”反转。
\end{enumerate}
\prosection{七、正方反驳库（针对反方全部可能论点）}
\begin{rebuttals}
\rebuttal{1}{同一份报告，91\% 是意愿题、46\% 是行为题；说的是为它，做的跟着我的钱走}{37\%、37\%、44\% 也是行为题，让的是自己；行为题两边都不过半}{打论据。反方要评委只看行为，那就只看行为：三成七已经削了自己的伙食水电，三成七已经真的负债，四成四差点放弃但付了，全是“过去一年”“已经”。反方拿到的是“有人让它等过”，没有拿到“多数让它等”。再看方向：一般就医拖过的四成六，急诊说会拖的只有一成六，反方自己也说假设题给方向，方向在正方。意图—行为鸿沟只说明意愿大于行为，没说行为反向}{追问：“付不起就让它等，这就是让路。” → 接受。付不起的人一部分借钱削口粮，一部分让它等，人数相当；反方要的“更”在哪里？}{最终}
\rebuttal{2}{病老终点它没有发言权，决定跟着我的负担和不舍走（依恋主导 69.3\%、结构化工具 7.3\%、负担中介显著）}{它的疼通过我的负担到达我的决定，渠道是我，起点和理由是它}{打机制。反方引的中介路径第一箭就是“它的疼痛与功能受损 → 我的负担”，我的负担是它的病的函数；“跟着负担走”翻译过来是“跟着它的病落在我身上的量走”。真正做出的决定：8,955 名犬主 87\%、646 名犬主 48.5\% 理由是它的疼，疾病与衰退合计 94.9\%，依据观察它 53.5\%、听兽医 39.2\%。69.3\% 不带符号：依恋深的人可能拖，也可能不忍它再疼而早放手。“它的指标是我读出来的”：观察的对象是它，兽医看的也是它；反方要证“感知随依恋变”，没有这个数据}{追问：“同样的疼，负担不同，决定不同。” → 那是决定有多难，不是决定的内容；反方需要“病一样、负担不同、决定不同”的数据，E16 是横截面，没有做过这个比较。追问：“Zarit 量表是为照护失智父母的子女设计的。” → 对，按反方的读法他们也“更爱自己”，一把判母爱也更是我的尺子分不出宠物之爱}{最终}
\rebuttal{3}{它的处境跟着我的状况走，我的状况不随它的处境波动，这个不对称就是让路方向}{反方引的每一份研究测的都是我的状况随它的处境波动}{打论据。238 人对照：它病了，我的负担、抑郁、生活质量全变，p 小于千分之一；157 人：它瘫了，我的负担显著升高；277 人：它的疼到我的负担，显著。不对称不存在，关系是对称的，对称只能看谁让路，病老期让路的是我的睡眠、收入、健康。“决定变量在我”是被反方自己淘汰的受益结构判准换了名字}{追问：“代价落在我身上不等于我跟着它决定。” → 判准两问：让谁受益、让谁让路。代价落在我身上，就是我让了路；受益的是被照护的它}{最终}
\rebuttal{4}{知道了也不改：喂胖、穿衣、当孩子，它的身体为我的方式让路（61\% 无应对、58\% 否认它喘）}{翻译错了，不等于翻译给自己}{打机制与论据。喂它是因为以为它想吃，弄错它的需要仍是为它做事；这是“怎么爱”不是“更爱谁”。61\% 那份问卷问的是“了解肥胖危害但目前无应对”，分母是全体饲主，而同一份报告说超重只有两到三成，七八成的宠物不胖，主人当然无应对；它不是“被告知它胖了却不改”的比例。58\% 是 31 例，反方自己承认是案例。给它穿衣可致热休克，只有“可致”没有发生率}{追问：“被告知之后呢？” → 被告知后仍不改，正方承认这算为我，请反方给这个比例。反方拿不出。追问：“被告知孩子超重仍喂零食的母亲更爱自己吗？” → 正方正要问这一句}{最终}
\rebuttal{5}{养宠是为了缓解孤独、被需要，起点就是我的需要（拟人化三因素、自我延伸）}{起点是需要，不等于指向是自己；受益结构标准下母爱也更是我}{打判准。所有关系的起点都是需要，交朋友因为需要朋友；“起点是我的需要”若等于“爱的是自己”，人类就没有一种爱指向别人。青年志 589 人里养宠首因是“一直喜欢小动物”78.3\%，“满足情感需要”36\%。1,359 人研究：依恋越高抑郁孤独越高，若为缓解孤独，依恋应减少孤独，数据相反，依恋在让我付出}{追问：“你承认得到陪伴就是承认为我。” → “更”是比较，回报存在不等于回报更重；回报归零时投入还在，才是比较的答案}{淘汰}
\rebuttal{6}{弃养：生活一变，被移走的是它（月饼、懒懒、流浪动物）}{弃养上新闻恰因少见；反方自己承认不比弃养人数、月饼只证机制}{打样本。救助站在结果上取样，看不见带着狗换城市的人；2020 年上海首例弃养案罚款 500 元，说明弃养是被处罚的例外；VetCompass 近三万例确认死亡，89.3\% 在诊所以安乐死结束，中位年龄 12.1 岁，典型的狗被养到老。“流浪动物 4000 万”出处未标，反方自己标了待核实}{追问：“阿派救助的老年宠物占十分之一。” → 分母是救助站收进来的动物，不是养老年宠物的家庭；它只说明被送走的老年宠物在救助里是少数}{淘汰}
\rebuttal{7}{按审美挑品种：明知短头犬会喘还买，选择每天在重复}{挑品种是买之前的事；进家后正因为爱的是这一只，明知会喘还养到老}{打时间范围。定义里“对宠物的爱”从进家算起。进家之后每天重复的是照顾它的选择，不是买它的选择；若为审美，换只不喘的更省事，没有人把它退回去。58\% 否认它喘是认知，31 例，是案例}{追问：“进家后 58\% 还在否认。” → 认知缺失不是相撞：不知道有冲突，何来让谁让路}{淘汰}
\rebuttal{8}{安乐死拖延与无效治疗：舍不得，拖着让它多疼（79\%、99.2\% 兽医被要求过）}{拖的是我的不舍，签字的是它的痛苦；兽医比例不是饲主比例}{打论据。反方已承认 79\%、99.2\% 是兽医一辈子被要求过的比例；饲主端：89.3\% 安乐死、8.3\% 自然死，拖到死的是少数；理由 87\%、48.5\% 是它的疼。拖延不带来享受，带来的是每天先看它还在不在的折磨；最后签下那张单子的是同一个人，在自己完全不愿意的情况下为了让它少疼，亲手结束了自己还能被它陪着的每一天}{追问：“自然死亡 8.3\% 也可能是拖死的。” → 那也只有 8.3\%。追问：“早放手也可能是负担高。” → 反方自己承认按它的指标放手是为它，不能两头都算}{最终}
\rebuttal{9}{葬礼、克隆安慰的是活人；哀伤 21\% 一生最痛，痛的是我}{离世之后出了时间范围；克隆上新闻正因几乎没人做}{打贴合度。反方自己已把克隆移出论点；PGD 7.5\%、21\% 最痛说明多数人接受它不可替代，克隆是替代，恰恰是少数人做的事；57.1\% 觉得必须藏起哀伤，若仪式是为我，人不会藏。“更是”预设两边都在，正方从不否认我会痛}{追问：“人为亲人办葬礼也是为生者。” → 是，人类学一直这么理解，这一点不丢人，但它不能当“更是我”的证据，哀伤阶段没有冲突可检验}{常见}
\rebuttal{10}{它老了怎么办（71.6\%）：担忧的是我的未来、我的负担}{担忧它的老年，正说明爱它；反方自己淘汰这一条时就是这句}{打机制。安排的是谁的老年？主体是我，内容是给它养老，受益是它。担心它老了怎么办的人，正是后来照护到 p 小于千分之一的那批人；美国两千人里 71\% 说它的未来一直在自己的担忧里}{追问：“困惑的主语是我怎么办。” → 判准看受益与让路，不看谁在动脑子}{淘汰}
\rebuttal{11}{宠物是可控的关系，它不会离开、不会反驳}{病老死是最不可控的阶段，投入却最高}{打事实。它失禁、痴呆、随时会走，我控制的只是自己要不要起床；若为可控，失控时应退出，多数没有；71.6\% 的年轻人困扰“它老了怎么办”，一段让人提前十年焦虑的关系不是可控的}{追问：“安乐死是终极的控制。” → 决定的内容是它的痛苦不是我的方便，89.3\% 按它的疼决定}{常见}
\end{rebuttals}

\consection{八、反方反驳库（针对正方全部可能论点）}
\begin{rebuttals}
\rebuttal{1}{让路的是我：九成愿负债、三成七削口粮、六成五先砍自己}{同一份报告让自己的三成七、让它等过的四成六，人数相当；过半的数字全是嘴上说的}{打论据。意图—行为鸿沟：说会负债的九成一，真负债的三成七；“先拖后付”先让路的是它；16\%（假设题）与 46\%（全年行为）不能连成曲线；91\% 是宠物医疗融资公司委托问出来的。行为题打平之后看它等不等跟着什么走：跟着我的余额}{追问：“三成七削口粮不算让自己？” → 算。我方不否认你爱它。行为题平局之后看决定变量：它等不等，跟着我的余额，不跟着它的病情}{最终}
\rebuttal{2}{回报归零，投入最高（238 人对照 p<0.001、瘫痪犬主人）}{被需要、可控感、自我形象写在共同定义里，病老期是它们的峰值；负担量表测的是代价，不是投入曲线}{打机制。代价我方承认，代价证明爱是真的，不证明爱向着谁；Zarit 量表测“照护对我自己生活的影响”；1,359 人研究依恋越高抑郁越高，抑郁跟着我的依恋走；同一作者 2022 年：负担越高越考虑安乐死。正方自己说“它的疼通过我的负担到达我的决定，渠道是我”，渠道是我，决定就跟着渠道的宽窄走}{追问：“被需要让人抑郁，谁会为自己追求抑郁？” → 不离婚、不辞职都会抑郁；依恋本身预测抑郁；抑郁说明代价，不说明方向}{最终}
\rebuttal{3}{终点按它来定：安乐死理由疼痛 48.5\%、87\%，疾病与衰退 94.9\%}{理由是事后填的，变量是当时起作用的；决策主导因素是我的依恋 69.3\%，用它的指标的只有 7.3\%}{打论据。疾病加衰老 94.9\% 对“让谁让路”无区分力，拖了两个月的家庭和当天决定的家庭填的都是“疾病”；自己观察 53.5\%，正方承认“所有疼痛都是被感知的”，那开关就在我对它的感知上；感知由依恋（69.3\%）、负担（中介显著）塑造}{追问：“69.3\% 没方向。” → 我方不用它定方向，用它定变量；方向由无效治疗（舍不得多治一段，79\% 兽医被要求过）给}{最终}
\rebuttal{4}{明知结局仍然开始，没有人为爱自己选剧痛（21\% 一生最痛、71.6\% 困惑它老了怎么办）}{85\% 说它是主要快乐来源；十几年最大的快乐换最后几个月的痛，任何为自己打算的人都会签}{打机制。结婚、生子同理；71.6\% 困惑“它老了怎么办”，困惑的主语是“我怎么办”；21\% 最痛的是我的失去；正方已承认这句是价值不是证据}{追问：“签账的人在快乐归零时该止损，可他们没有。” → 快乐归零时被需要与可控感到峰值；真止损的人不在照护样本里（抽样条件是正在照护）；46\% 拖过就医}{最终}
\rebuttal{5}{武汉封城 5000 只留守宠物：主人拿不到陪伴还要救它}{极端个案，双方约定只证机制；它们为什么在家？出门时“带不走就不带”是主人的日常决定}{打代表性。冒险进门的是 60 名志愿者、20 位开锁师傅，主人让出的是隐私；5000 只对 1.26 亿只；封城造成的是回不去，留下是出行前决定的。我方不指责，只说证不了比例}{追问：“封城是突然的。” → 对，所以这个案例双方都只说机制}{淘汰}
\rebuttal{6}{哀伤与葬礼：57.1\% 必须隐藏哀伤、PGD 7.5\%、宠物殡葬}{痛的是我，隐藏的是我的哀伤；葬礼安慰的是活人，它已经没有需要了}{打判准。冲突已结束，无从检验（正方备赛时自己因此淘汰这一条）；57.1\% 求的是社会承认我的哀伤；死后一切支出受益人只能是我，包包克隆 30 万只作画面}{追问：“人为亲人办葬礼也是为自己？” → 是为生者，这不丢人，我方不做道德指控；辩题问“更是”，哀伤阶段没有它的一份}{淘汰}
\rebuttal{7}{意志性关切：爱是意志的（Frankfurt）、它不可替代（Badhwar），所以爱指向对方}{这是把结论写进定义，双方约定不做定义杀；非替代性双方承认，不判方向}{打学理边界。稳健关切说描述理想之爱，辩题问典型饲主；“为它之故”的关切先要知道它的需要，84\% 狗主认为超重犬体型健康、58\% 短头犬主否认它喘；配偶也不可替代，没有人因此说丧偶之痛是为了配偶}{追问：“承诺不看回报。” → 承诺是我给自己立的，守承诺守的是“我是什么样的人”，共同定义里叫自我形象}{淘汰}
\rebuttal{8}{为它改了作息、辞了职、凌晨起床（日常让路）}{习惯就是我的生活，没有相撞就没有让路；正方自己因无可核实比例淘汰了这一条}{打论据。慧慧辞职是个案，上了深度报道恰说明不典型；月饼“要去外地”被送走，个案对个案；日常内容全由我定：几点遛、吃什么、穿不穿}{追问：“凌晨三点起床不是让路？” → 是，我方不否认。但起不起床由我定，它叫了我不起它就等，决定权在我，这正是可控感}{淘汰}
\rebuttal{9}{为它花了很多钱：3126 亿、单犬 3006 元、医疗占 27.6\%}{花多少说明爱有多重，花在哪说明爱向着谁}{打贴合度。同一份白皮书食品占 53.7\%，喂食最容易喂坏它；青年志服装配饰超五成、写真超两成、保健品超七成；美国样本里 46\% 因成本延迟就医、25\% 换便宜粮。不同样本只作画面对照，不相减}{追问：“写真是纪念它。” → 照片挂在我墙上，它不看照片；这不是坏事，但受益人是我}{常见}
\rebuttal{10}{它认得我、它爱我（陌生情境实验、皮质醇相关）}{它爱不爱我是它的爱的方向，辩题问的是我的爱的方向}{打贴合度。因变量全在动物身上；它证明得越漂亮，越说明我们拿到了一个真正会回应我们的对象；狗的皮质醇跟着主人的皮质醇走，它的状态跟着我的状态走，正是反方命题}{追问：“双向的爱当然更是它。” → 双向不改变各自的方向；它的爱我方不争，只争我的}{淘汰}
\rebuttal{11}{救助流浪动物、领养代替购买：不挑对象的爱}{救助是义举，不是辩题里的“对宠物的爱”；进家之后钱、病老、终点三个冲突一个不少}{打定义范围。共同定义从它进家门成为“这一只”开始算；领养时也在挑品种、年龄、性格；救助站里的懒懒、月饼是被让路的那一头；救助带来的道德满足感是真实的回报}{追问：“领养病猫老狗的人图什么？” → 不否认，个案；判准以多数典型为准}{常见}
\end{rebuttals}
\section{九、持方优劣评估}
\begin{assessment}
\assess{定义空间}{「爱」的中立定义双方都接受，不争。「更」的中立读法（冲突时让谁让路）是正方一辩稿钉下的，反方上一版最想用的正是这把尺子，如今由正方先拿在手里；反方要在同一把尺子下打，只能争“行为为准、多数为准”，空间比正方小。}{偏正}
\assess{论证义务}{中立判准写着“以多数饲主典型表现为准”。钱这个战场上，两方的行为题都不过半（削自己 37\%、让它等过 46\%），谁也证不了多数；正方在病老与终点战场有过半的数字（安乐死理由是它的疼 87\%、疾病与衰退 94.9\%），反方在决定变量上有过半的数字（依恋主导 69.3\%、凭自己观察 53.5\%）。两边义务都重，正方略轻，因为“投入不降反升”只需要证明照护者仍在照护。}{略偏正}
\assess{论据可得性}{可核实的论据两边都够：正方有负债与削口粮、照护负担 p<0.001、终点理由三组；反方有同一份报告的 46\%、同一作者的负担预测安乐死考虑、依恋主导 69.3\%、无效治疗 79\%。反方的硬伤是“让路”这个动作没有一条过半的直接行为数据，中国数据只有 61\% 且分母不对；正方的硬伤是多数比例全在英美，中国只有账单量级与两个案例。}{均衡}
\assess{评委直觉}{只用一般性判断：新生赛评委多为高年级辩手与教练，偏好定义清楚、正面回应、不玩定义杀。养宠的人听“你爱的其实是你自己”会本能反感，反方要先克服直觉；但评委也听得出“付钱是为它、停钱也是为它、撑着也是为它”的不可证伪味道，正方必须一开场就给出“什么算为我”的三条，否则直觉优势会被质询吃掉。}{偏正}
\assess{赛制顺序}{正方先立论、最后结辩：正一必须在 20 秒内把判准说清，因为反四第一条链就冲“九成是意愿题”。反方先结辩：反四要封住正四“回报归零仍在喂药”的画面。正四留 30–40 秒真回应。结构上均衡。}{均衡}
\end{assessment}
\begin{fields}
\claim[结论]{\textbf{反方（爱的更是我）较劣势}。劣势来自定义空间与评委直觉：判准被正方先定在“冲突时优先谁”，而“让路的是它”这个动作反方拿不出一条过半的行为数据。反方的补偿是两处同源反证（同一份报告的 46\%、同一作者的 2022 年中介效应）和一条过半的决定变量数据（依恋主导 69.3\%）。}
\begin{concard}{劣势方打法（反方）}\textbf{收窄战场}：不打“它让路的多数”，打“决定变量在我”：钱的战场只求打平（行为题两边不过半），病老终点战场用 E16、Kiss 69.3\%、7.3\% 立“跟着我的状况走”。\textbf{抢先定义}：反一开头就说“同一份报告，正方引意愿题，我方引行为题”，把评委的权重先定在行为上。\textbf{判准之争}：不换判准，只加一句“以行为为准”。\textbf{价值层}：收在“承认之后才能真的看它一眼”，绝不收在犬儒；每个环节先说“我方不否认你爱它”。\textbf{顺序利用}：反四封路式，提前点破“老猫喂药”的画面里“没被说出来的那一半”。\end{concard}
\begin{procard}{优势方要防的（正方）}三处最容易翻车：（1）被问“被需要也是你的需要吗”答“不是”，一答就被追到不可证伪；口径是“需求落在我身上，不等于回报落在我身上”。（2）在 E16 上说出“撑到撑不住才停”，那是把停止开关交给反方；口径是“开关在它的指标上，我的负担只决定我有多痛苦”。（3）用 84\%（假设题）证多数，被反方指出和 91\% 同一题型；口径是“行为题两边不过半，方向看三层”。另外三条铁律：不说“我什么都没得到”；不说“所以它也爱我”；不拿“我为它花了多少钱”当爱它的证据。\end{procard}
\end{fields}
\section{十、口径表}
\prosubsection{10.1 正方口径表}
\begin{canon}{正方}
\canongroup{定义}
\canonrow{爱}{对一个对象的深挚感情，以及这份感情带来的持续关切和行动}{有感受，也有为它做的事}{“爱就是为对方之故”（定义杀）}{汉典；SEP}
\canonrow{更}{这是一道比较题，两边都在。我方承认养宠让人快乐、让人被需要，我方争的是分量}{比的是权重，不是有无}{“这份爱里没有半点自我满足”“我什么都没得到”}{}
\canonrow{它}{你家那一只，它有自己的需要，不随你的想法改变}{会区分人的真实他者}{“所以它也爱我”（一说就被反二接走）}{}
\canongroup{判准}
\canonrow{冲突时优先谁}{当它的需要和我的需要撞上，让路的是谁，这份爱就更是谁的；说全两问：让谁受益、让谁让路}{冲突检验；有得选的时刻谁让路}{只说“让谁让路”不说“受益”；“回报归零时投入是否还在”当唯一判准}{备赛文档第四节}
\canonrow{双方义务}{我方要证明多数人让的是自己；对方要证明多数人让的是它}{在真实冲突里典型饲主让的是谁}{“我方只需证明有人让自己”}{}
\canongroup{论点}
\canonrow{回报归零，投入最高}{它病了、老了、不再回应，它能给我的一切归零；主人的投入不但没有减少，反而在自己身上留下了可测量的伤害。图回报，解释不了这件事}{回报最低的时候，投入最高；投入就是让出去的睡眠、时间、钱、健康}{“被需要不是我的需要”“我的负担不决定它哪天走”“撑到撑不住才停”}{E1、E17、E29、E23}
\canonrow{让路的是我}{削减自己的口粮去付它的药费，这就是“冲突时优先谁”最朴素的答案}{钱和它的病相撞时，让路的是我}{“行为题里多数让自己”“八成四证明多数”}{E3、E4、E20}
\canonrow{明知结局，仍然开始}{没有人会为了爱自己，去主动选择一段注定在十几年后带来剧痛的关系。我们明知结局，仍然开始}{终点按它来定：多数人按它的痛、不按自己的不舍做决定}{把这句当证据用（它是价值收束）}{E5、E24、E27、E26}
\canonrow{全队第一句}{它的疼通过我的负担到达我的决定；渠道是我，起点和理由是它}{负担影响决定有多难，不改变决定的内容是它的痛}{“负担与决定无关”}{E16、E28、E24}
\canonrow{全队第二句}{被需要是我的需要，我方承认；它需要我起床、我需要睡觉，起床的是我}{被需要得到满足，是因为我让了路}{“被需要不是回报”}{}
\canonrow{什么算为我}{三样：有得选时让它扛；拖着不让它走让它多疼一段；被告知它有问题之后仍不改。都存在，典型饲主做的是相反的事}{可证伪标准}{“任何行为都是为它”}{}
\canongroup{数据}
\canonrow{照护负担}{二零一七年一项对照研究比较了二百三十八位猫狗主人，其中一百一十九位正在照护患慢性或终末期疾病的动物；照护组负担、压力、抑郁与焦虑症状全面更高，生活质量更差，每一项差异统计上都极其显著（p<0.001）；研究者的概括是，这种损耗在许多方面与照护一位患痴呆的家人相似}{Spitznagel 2017，Veterinary Record；“痴呆”是作者新闻稿原话}{报“p 值小于零点零零一”这类小数；说“论文里写了痴呆”}{E1、E2}
\canonrow{负债与削口粮}{二零二六年一项两千人的调查中，九成主人表示愿意为救宠物负债；过去一年有三成七的人为支付兽医费削减了自己的食品、水电和其他账单}{九成一愿负债；三成七已负债；四成四差点放弃但付了；同一份报告四成六延迟或跳过就医、一成六遇两千美元急诊会拖}{“多数人削了自己”；用八成四证多数}{E3}
\canonrow{先砍自己}{美国两千名犬猫主里六成五会先削减自己的生活预算再动宠物的；七成一说它的未来一直在自己的担忧里}{意愿题，只说方向}{当行为题用}{E4}
\canonrow{终点理由}{六百四十六名犬主里安乐死首要理由是疼痛与痛苦，接近一半；八千九百多名犬主里八成七把疼痛与痛苦列为理由；十七国两百二十八名宠物主里理由是疾病六成三、衰退三成二，依据是观察它五成三、听兽医近四成}{48.5\%；87\%；63.3\%、31.6\%、53.5\%、39.2\%；同一摘要“依恋主导 69.3\%”不带符号}{把 69.3\% 改写成“决定有多难”}{E24、E27、E26}
\canonrow{死亡方式}{英国近三万例确认死亡里，近九成的狗在诊所以安乐死结束，中位年龄十二岁多}{89.3\%、8.3\%、12.1 岁}{用它证弃养比例}{E23}
\canonrow{哀伤}{英国九百七十五人里，经历过宠物死亡的人中两成一把它列为一生最痛的失去；延长哀伤障碍条件发生率百分之七点五}{21.0\%、7.5\%}{“所以爱的更是它”（非替代性双方承认，不单独判方向）}{E5}
\canonrow{它老了怎么办}{中国青年饲主里七成一困惑“它老了怎么办”}{71.6\%，只说“知道结局”}{说担忧的内容}{E13}
\canonrow{账单量级}{中国单只犬年均消费三千多元，医疗占近三成；狗跟腱手术五千五、猫 CT 八千}{3006 元、27.6\%、5500 元、8000 元}{“花了多少钱证明爱它”}{E7、E38}
\canongroup{案例}
\canonrow{卢卡}{新京报报道过一位女孩，她十五岁的哈士奇从十三岁起失禁、瘫痪、痴呆；她一晚上起床六七次，第二天照常上班，后来辞了职，每三小时回家一趟；狗十四岁多不能吞咽时，她接受了安乐死}{为它辞职，为它起床，最后也为它放手}{加原文没有的细节（“和兽医一起判断”）}{E21}
\canonrow{武汉封城}{主人把自己家的钥匙和身份证视频交给陌生人，让志愿者开锁进门，为将近五千只留守宠物补充食水}{让出隐私换它一口水；只说机制}{“冒风险的是主人”“证明多数”}{E14}
\end{canon}
\consubsection{10.2 反方口径表}
\begin{canon}{反方}
\canongroup{定义}
\canonrow{爱}{对一个对象的深挚感情，以及这份感情带来的持续关切和行动（与正方相同，不争）}{}{“爱就是满足自我需要”（定义杀）}{汉典；SEP}
\canonrow{更}{这是一道比较题，两边都在。我方不否认你爱它，我方比的是分量}{争的是重心，不是有无}{“人根本不爱动物”“养宠的人自私”“别骗自己了”}{}
\canonrow{我}{饲主自己的需要：陪伴、被需要、情感出口、可控感、自我形象}{不是道德指控}{“自私”}{}
\canongroup{判准}
\canonrow{冲突时优先谁}{当它的需要和我的需要撞上，让路的是谁；以行为为准，以多数为准}{同一把尺子，我方只加一句“以行为为准”}{换成“受益结构”“内容由谁塑造”当判准}{备赛文档第四节}
\canonrow{一句话命题}{在钱、病老、终点三个冲突时刻，让不让路跟着的是我的钱、我的负担、我的不舍，不是它的指标}{它的处境跟着我的状况走}{“我的状况不随它的处境波动”（被自引数据证伪）}{}
\canonrow{什么算为它}{三样：有得选时让自己扛；按它的指标而不是我的状况定终点；被明确告知后按它的需要改。都存在；有比例数据的地方，跟着我走的是多数}{可证伪标准}{“任何行为都是为我”“安乐死是它让路、继续治也是它让路”}{}
\canongroup{论点}
\canonrow{说的是为它，做的跟着我的钱走}{同一份报告，正方引的是意愿题，我方引的是行为题；问卷上愿意负债的是九成，账单真来时让它等过的近半}{它等不等，由我的余额决定}{“46 大于 37”“三十个百分点的鸿沟”“多数人让它等”}{E3、E29}
\canonrow{病老终点，舍不得就让它多受一段}{病老之时它一句话说不了；决定它多疼一段还是早走一步的，是我的负担和我的不舍；不舍的结果是它多受一段}{决定变量在我}{“早放手也是它让路”（两头都算）}{E26、E16、E30、E32}
\canonrow{对正方“回报归零”}{负担量表测的是照护对我自己生活的影响；负担高说明这段关系在我生命里有多重；同一作者说负担越高越考虑安乐死}{代价证明爱是真的，不证明爱向着谁}{“他们不爱它”}{E16、E40}
\canonrow{对正方“明知结局”}{八成五说宠物是主要快乐来源；十几年最大的快乐换最后几个月的痛，是任何为自己打算的人都会签的账}{正方已承认这句是价值不是证据}{}{E4}
\canongroup{数据}
\canonrow{行为题}{同一份两千人报告：四成六过去一年因成本延迟或跳过就医，四成四上次差点因钱放弃，一成六遇两千美元急诊会拖；另一份一千人调查两成五改买更便宜的粮}{46\%、44\%、16\%、25\%；承认 37\%、37\% 也是行为题}{“46 大于 37”}{E3、E20}
\canonrow{意图—行为鸿沟}{心理学的元分析：意愿的大幅变化只带来行为的小幅变化}{d+ = .36，Sheeran \& Webb 2016}{把 16\% 与 46\% 相减}{E29}
\canonrow{决定变量}{十七国两百二十八名宠物主：情感依恋是安乐死决策的主导因素，近七成；依据是自己观察五成三，用结构化生活质量工具的只有百分之七}{69.3\%、53.5\%、7.3\%；同一摘要理由疾病 63.3\%、衰退 31.6\%，正方会引}{说 69.3\% 给出了方向}{E26}
\canonrow{负担中介}{两百七十七名骨关节炎犬主人：照护负担对“考虑安乐死”有显著间接效应}{同样的疼，负担不同，决定不同}{“负担高就安乐死”（是“考虑”）}{E16}
\canonrow{无效治疗}{八百八十九名北美兽医里近八成被要求过提供自己认为无效的治疗；四百七十七名兽医里八成五过去一年遇到过；研究者在评论文中说七成六的兽医认为无效治疗有时有益的是主人}{79\%、85.0\%、76.2\%；兽医比例不是饲主比例}{当饲主比例用}{E30、E31、E32}
\canonrow{内疚}{英国一项调查：过半主人在宠物死后感到内疚，确信时机正确的仍有近半内疚，四分之一后悔}{53\%、47\%、25\%；样本量未公布}{报样本量}{E36}
\canonrow{依恋与抑郁}{一千三百多名英国饲主：依恋越高，抑郁、焦虑、孤独越高}{1,359 人，Ellis 2024}{“所以养宠有害”}{E11}
\canonrow{日常追问}{中国一千三百多份问卷：被提示宠物肥胖后六成一表示了解但没有应对；短头犬主人近六成否认它有呼吸问题；给它穿衣可致热休克甚至死亡}{61\%（分母是全体饲主，只用来追问）；58\%（31 例，案例）}{用 61\% 证多数；“不到一小时中暑”}{E33、E8、E10}
\canonrow{账单量级}{单只犬年均三千多元；狗跟腱手术五千五、猫 CT 八千}{3006 元、5500 元、8000 元}{}{E7、E38}
\canongroup{案例}
\canonrow{月饼}{中秋节那天送到北京房山救助站的八哥犬，它的主人说自己要去外地，家里其他人都不想继续养}{生活变动时被移走的是它；只证机制}{“证明多数人弃养”}{E15}
\canonrow{卢卡（反用）}{正方那只哈士奇从十三岁失禁瘫痪痴呆到十四岁多不能吞咽，中间将近两年；开关是“还能不能活”，不是“过得好不好”}{不评价对错，只问开关跟着什么走}{“她害了它”}{E21}
\canonrow{包包克隆}{一只德牧安乐死之后，主人花三十万以上克隆它；它已经不在了，三十万买的是我的不失去}{只作结辩画面，不作证据（离世之后）}{当论据用}{E37}
\end{canon}
\section{十一、论据出处清单}
\begin{sourcelist}
\source{1}{238 位猫狗饲主对照研究：119 位照护慢性或终末期疾病动物 vs 119 位按年龄、性别、物种匹配的健康对照；照护组负担、压力、抑郁与焦虑症状更高，生活质量更差，全部 p<0.001；负担越高心理社会功能越差}{Spitznagel MB, Jacobson DM, Cox MD, Carlson MD (2017). Caregiver burden in owners of a sick companion animal: a cross-sectional observational study. Veterinary Record 181(12):321. DOI 10.1136/vr.104295（摘要经 Europe PMC 核实）}{greater burden, stress and symptoms of depression/anxiety, as well as poorer quality of life (p<0.001)}{2026-09-18}{已核实}{正方论点二}{}
\source{2}{作者原话：照护病宠的损耗“在许多方面与照护患痴呆的家人相似”；600 名宠物主在线问卷}{Kent State University 新闻稿《When Caring for a Sick Pet Becomes Too Much》2017-09-18，\url{https://www.kent.edu/kent/news/when-caring-sick-pet-becomes-too-much}}{in many ways similar to what we see in a person caring for a sick family member, for example, a parent with dementia}{2026-09-18}{已核实}{正方论点二}{新闻稿转述，不是论文原文；场上说“研究者的概括”}
\source{3}{2,000 名美国宠物主（Censuswide，2026-05-21 至 27）：91\% 愿为救宠物负债；37\% 过去一年用信贷或负债付兽医费；37\% 削减食品杂货、水电、账单；46\% 因成本延迟或跳过就医；44\% 上次就医差点因成本放弃；16\% 遇 2,000 美元急诊会延迟或放弃}{Lovet《2026 Checkup Report》，PR Newswire 2026-06-24，\url{https://www.prnewswire.com/news-releases/91-of-pet-owners-would-go-into-debt-to-save-their-pets-life-yet-nearly-half-have-delayed-care-due-to-cost-new-report-finds-302808589.html}；lovet.com 博客 2026-06-30}{91\% of pet owners say they would be willing to take on debt to save their pet's life; 37\% have cut back on everyday expenses such as groceries, utilities, or bills due to veterinary costs; 46\% of pet owners reported delaying or skipping veterinary care due to rising costs}{2026-09-18}{已核实}{双方（正方论点一、反方论点一）}{}
\source{4}{2,000 名美国犬猫主（Talker Research 为 MetLife 宠物保险做，2025-12-12 至 19）：65\% 会先削减自己的生活方式预算再动宠物的；85\% 说宠物是 2025 年主要快乐来源；91\% 说宠物给了正向关注点；71\% 说宠物的未来一直在自己的担忧里}{Talker Research《Most pet owners would cut personal spending before sacrificing pets' needs》2026-01-30，\url{https://talkerresearch.com/most-pet-owners-would-cut-personal-spending-before-sacrificing-pets-needs/}}{Sixty-five percent of pet owners would make budget cuts to their own lifestyle before disrupting their pet's.; 85\% said their pet has been their main source of happiness throughout 2025.}{2026-09-18}{已核实}{双方}{}
\source{5}{英国 975 人全国配额样本：32.6\% 经历过宠物死亡，其中 21.0\% 把宠物之死列为一生最痛的失去；宠物死亡后延长哀伤障碍条件发生率 7.5\%；宠物之死占全部延长哀伤障碍病例 8.1\%}{Hyland P (2026). No pets allowed: Evidence that prolonged grief disorder can occur following the death of a pet. PLOS ONE. DOI 10.1371/journal.pone.0339213}{21.0\% of these people chose the death of their pet as most distressing; The conditional rate of PGD following the death of a pet was 7.5\%}{2026-09-18}{已核实}{正方论点三、反方论点二}{}
\source{6}{RSPCA 2026 宠物哀伤调查：57.1\% 觉得必须隐藏哀伤；13.4\% 觉得哀伤被完全理解；6.9\% 认为宠物哀伤被足够认真对待；0.7\% 把宠物只当“宠物”}{RSPCA《What our survey data tells us about pet grief》（Pet Grief Survey 2026），\url{https://www.rspca.org.uk/adviceandwelfare/pets/bereavement/survey}}{57.1\% of people felt they had to hide their grief.}{2026-09-18}{已核实}{双方备用}{慈善机构自愿参与样本}
\source{7}{2025 年城镇犬猫 1.26 亿只；消费 3126 亿元；医疗 27.6\%、食品 53.7\%；单只犬年均 3006 元、猫 2085 元}{《2026 年中国宠物行业白皮书（消费报告）》派读宠物行业大数据出品，央视网 2026-01-05，\url{https://business.cctv.com/2026/01/05/ARTIdmv2AWmVeRVodaZYLOvQ260105.shtml}}{2025年城镇犬猫数量为1.26亿只；医疗市场，市场份额为27.6\%}{2026-09-18}{已核实}{双方}{}
\source{8}{285 只犬、68 品种、31 例短头阻塞性气道综合征；受影响犬的主人 58\% 说狗没有呼吸问题}{Packer RMA, Hendricks A, Burn CC (2012). Animal Welfare 21(S1):81–93. DOI 10.7120/096272812X13345905673809（Cambridge Core 摘要）}{Over half (58\%) of owners of affected dogs reported that their dog did not have a breathing problem.}{2026-09-18}{已核实}{反方日常战场}{n=31，只作案例}
\source{9}{59\% 的狗、61\% 的猫超重或肥胖（2022 流行率调查）；84\% 狗主人、70\% 猫主人认为体型健康（2023 调查，527 名美国宠物主）}{Association for Pet Obesity Prevention，\url{https://www.petobesityprevention.org/2023}}{61\% of cats and 59\% of dogs are overweight or have obesity; 84\% of dogs and 70\% of cat owners assessing their pets' body condition as healthy}{2026-09-18}{已核实}{反方日常战场}{}
\source{10}{拟人化行为干扰体温调节、脱水、热休克甚至死亡；垃圾食品与热量失衡导致肥胖}{Mota-Rojas D 等 (2021). Anthropomorphism and Its Adverse Effects on the Distress and Welfare of Companion Animals. Animals 11(11):3263. DOI 10.3390/ani11113263（PMC8614365）}{interfering with thermoregulation ... heat shock, even death}{2026-09-18}{已核实}{反方日常战场}{“不到一小时中暑”这一具体时长未见于摘要，场上只说“可致中暑甚至死亡”}
\source{11}{1,359 名英国犬猫饲主：自我延伸正向预测积极情绪、负向预测孤独；依恋正向预测抑郁、焦虑、孤独}{Ellis A, Stanton SCE, Hawkins RD, Loughnan S (2024). Animals 14(3):441. DOI 10.3390/ani14030441（PMC10854534）}{attachment ... a significant positive predictor of depression, anxiety, and loneliness}{2026-09-18}{已核实}{反方论点二、正方反驳库}{}
\source{12}{29 对犬与主人的陌生情境实验（26 对入分析）：狗对主人更多接近与接触；狗与主人末次皮质醇正相关}{Ryan MG, Storey AE, Anderson RE, Walsh CJ (2019). Frontiers in Behavioral Neuroscience 13:162. DOI 10.3389/fnbeh.2019.00162}{Final CORT levels for dogs were positively correlated with their owners' final CORT levels.}{2026-09-18}{已核实}{正方反驳库（“真实他者”）}{}
\source{13}{青年志 648 份问卷、589 人养过宠：78.3\% 因“一直喜欢小动物”；36\% 满足情感需要；71.6\% 困惑“它老了怎么办”；现实中 41.7\% 自称父母、理想中 25.4\%；服装配饰超 50\%、写真超 20\%、保健品超 70\%}{青年志 Youthology《宠物本位时代——2025 青年养宠观报告》腾讯新闻 2025-02-28，\url{https://news.qq.com/rain/a/20250228A02VHN00}}{71.6\%的受访者对宠物‘老了怎么办’感到困惑，位列养宠困惑之首。}{2026-09-18}{已核实}{双方}{非概率、青年为主、样本小}
\source{14}{2020 年武汉封城：协会核心人员增至 8 人、60 名志愿者、约 20 名开锁师傅；1 月 26 日晚至 2 月 11 日为将近 5000 只宠物补充食水；求援者手持身份证录视频委托上门}{腾讯新闻谷雨工作室 2020-02-28《武汉 8 人协会如何 1 个月援救 5000 只留守宠物》，\url{https://news.qq.com/rain/a/20200228A07NDT00}}{一共为将近5000只宠物补充了食物和水}{2026-09-18}{已核实}{正方论点一案例}{极端情境，只说机制}
\source{15}{北京房山救助基地八哥犬“月饼”2024 年中秋当天送来，主人要去外地、家人不愿再养；2020 年上海首例遗弃犬只案罚款 500 元}{中国日报网转法治日报 孙天骄 2025-01-23《弃养宠物行为多发 如何提高养宠的责任门槛？》，\url{https://cnews.chinadaily.com.cn/a/202501/23/WS6791dc8da310be53ce3f315b.html}}{‘月饼’主人就联系我们，说自己要去外地，而家里其他人都不想继续养犬了。}{2026-09-18}{已核实}{反方论点一案例}{个案}
\source{16}{277 名骨关节炎犬主人（社交媒体招募）：疼痛与功能受损、主人考虑安乐死、照护负担、治疗满意度之间关系显著（P≤0.01）；照护负担对“考虑安乐死”有显著间接效应}{Spitznagel MB 等 (2022). Veterinary Journal 286:105868. DOI 10.1016/j.tvjl.2022.105868（摘要经 Europe PMC 核实）}{a statistically significant indirect effect of caregiver burden}{2026-09-18}{已核实}{反方论点二、正方论点二口径}{}
\source{17}{67 名瘫痪或轻瘫犬主人 vs 90 名健康犬主人：Zarit 负担显著更高（p<0.001）、生活质量更低（p=0.03）、压力更高（p=0.04）；抑郁焦虑无显著差异}{Pryjmak J, Zablotski Y, Lauer SK (2026). Frontiers in Veterinary Science. DOI 10.3389/fvets.2026.1837756}{Owners of paretic/paralyzed dogs reported significantly higher ZBI scores (p < 0.001).}{2026-09-18}{已核实}{正方论点二}{德国样本，88.5\% 女性}
\source{18}{Frankfurt 1999 p.129；Badhwar 2003 p.63；Nozick 1989；Whiting 1991 / Delaney 1996}{Stanford Encyclopedia of Philosophy「Love」条（Helm, B., 2021 修订版），\url{https://plato.stanford.edu/entries/love/}}{neither affective nor cognitive. It is volitional; phenomenologically non-fungible}{2026-09-18}{已核实}{双方定义与学理}{}
\source{19}{「爱」第一义项“对人或事有深挚的感情”}{汉典 zdic.net「爱」字条（《现代汉语词典》体系，2026 年访问），\url{https://www.zdic.net/hans/爱}}{对人或事有深挚的感情}{2026-09-15}{已核实}{双方定义}{}
\source{20}{Rover 2025（Pollfish，1,000 名美国宠物主，2025-02）：34\% 说宠物支出是预算紧张时最后削减的类别之一；33\% 已在其他生活领域减支以保证宠物所需；25\% 改买更便宜的宠物食品或服务}{Rover《True Cost of Pet Parenthood Report 2025》新闻稿 2025-03-18，GlobeNewswire，\url{https://www.globenewswire.com/news-release/2025/03/18/3044277/0/en/Rover-Releases-True-Cost-of-Pet-Parenthood-Report-for-2025.html}}{A third (34\%) of pet parents say pet spending is one of the last categories they would cut}{2026-09-18}{已核实}{双方}{}
\source{21}{慧慧 15 年前带回哈士奇卢卡；13 岁起失禁、器官衰竭、瘫痪、痴呆；一晚最多起床六七次第二天照常上班；2024 年辞职专心照顾，每 3 小时回家 1 次；卢卡 14.9 岁因无法吞咽接受安乐死；阿派救助的老年宠物占十分之一}{新京报 秦冰 2025-05-24《老年宠物，年轻人的爱与“不可承受之重”》，\url{https://m.bjnews.com.cn/detail/1748051658168890.html}}{2024年，慧慧辞职，专心照顾卢卡。；最高纪录是一晚上起床六七次，第二天再去正常上班。}{2026-09-18}{已核实}{正方论点二、三案例；反方反用}{}
\source{22}{英国 VetCompass 2016 年 905,544 只犬；随机抽样 29,865 例确认死亡，26,676（89.3\%）为安乐死、8.3\% 自然死亡；安乐死犬中位年龄 12.1 岁}{Pegram C 等 (2021). Proportion and risk factors for death by euthanasia in dogs in the UK. Scientific Reports 11:9145. DOI 10.1038/s41598-021-88342-0（PMC8096845）}{Of the confirmed deaths, 26,676 (89.3\%) were euthanased and 2,487 (8.3\%) died unassisted.}{2026-09-18}{已核实}{正方论点二、三}{只量记录在案的死亡，不用来说弃养比例}
\source{23}{Dog Aging Project 终末调查 646 名犬主：536（83.0\%）安乐死；首要理由疼痛或痛苦 48.5\%、生活质量差 24.8\%、预后差 19.6\%}{McNulty KE 等 (2026). JAVMA 264(7):892–901. DOI 10.2460/javma.25.12.0863（PMC13340034）}{The most common reason for euthanasia was pain and/or suffering (260 of 536 [48.5\%]).}{2026-09-18}{已核实}{正方论点三}{自报、志愿队列}
\source{24}{共情—利他假说“易逃离”实验：高共情组逃离容易时帮助率 .91、困难时 .82；低共情组 .18 与 .64（每格 11 人）}{Batson CD, Duncan BD, Ackerman P, Buckley T, Birch K (1981). Is Empathic Emotion a Source of Altruistic Motivation? Journal of Personality and Social Psychology 40(2):290–302，\url{https://greatergood.berkeley.edu/images/uploads/Baston-EmpathySourceAltruism.pdf}}{subjects feeling a high degree of empathy for the victim should be as ready to help when escape without helping was easy as when it was difficult}{2026-09-18}{已核实}{正方论点一学理}{人对人实验室研究，只用判别逻辑}
\source{25}{河南郑州吴先生养 8 年的宠物狗产后瘫痪，配轮椅带它散步，“要陪它到老”}{澎湃新闻 王文娟 2020-05-08《宠物狗瘫痪坐轮椅出行，主人：把它当孩子》，\url{https://www.thepaper.cn/newsDetail_forward_7295493}}{已经养了8年}{2026-09-18}{已核实}{正方备用案例}{短视频新闻，细节少}
\source{26}{228 名宠物主（17 国）与兽医访谈：安乐死理由疾病 63.3\%、年龄相关衰退 31.6\%；判断依据自己观察 53.5\%、兽医建议 39.2\%、结构化生活质量工具 7.3\%；情感依恋是决策主导因素 69.3\%}{Kiss A, Möller W, Wagenhoffer Z, Fodor K (2026). Factors Influencing Decision-Making in Companion Animal Euthanasia: A Mixed-Methods Study of Pet Owners and Veterinarians. Animals 16(11):1738. DOI 10.3390/ani16111738（摘要经 Europe PMC 核实，PubMed 42278170）}{Illness (63.3\%) and age-related decline (31.6\%) were the most frequently reported reasons; emotional attachment was the dominant factor in decision-making (69.3\%)}{2026-09-18}{已核实}{双方（同一摘要各引一半）}{}
\source{27}{Dog Aging Project 8,955 名犬主：7,764（87\%）把疼痛或痛苦列为安乐死理由；近期兽医就诊与饲主感知疼痛正相关}{Mann AM, O'Brien J, Dog Aging Project Consortium, Ruple A, Parker RL (2026). American Journal of Veterinary Research（摘要经 Europe PMC 核实），\url{https://pubmed.ncbi.nlm.nih.gov/42526497/}}{Overall, 7,764 of 8,955 owners (87\%) reported pain and/or suffering as a reason for euthanasia}{2026-09-18}{已核实}{正方论点三}{自变量是饲主感知的疼痛}
\source{28}{393 名年老或重病伴侣动物的饲主：照护负担、预期性哀伤、生活质量三者不可互换}{Spitznagel MB, Anderson JR, Marchitelli B, Sislak MD, Bibbo J, Carlson MD (2021). Veterinary Record 188(9):e74. DOI 10.1002/vetr.74（摘要经 Europe PMC 核实）}{caregiver burden, anticipatory grief and QoL are not interchangeable and may differentially impact owner decisions and behaviour.}{2026-09-18}{已核实}{正方论点二口径}{}
\source{29}{意图—行为鸿沟：操纵意愿的实验元分析，意愿中到大幅变化只带来行为小到中幅变化 d+ = .36}{Sheeran P, Webb TL (2016). The Intention–Behavior Gap. Social and Personality Psychology Compass 10(9):503–518. DOI 10.1111/spc3.12265（White Rose 作者手稿 \url{https://eprints.whiterose.ac.uk/id/eprint/107519/}）}{led to only a small-to-medium-sized change in behavior (d+ = .36; Webb \& Sheeran, 2006)}{2026-09-18}{已核实}{反方论点一学理}{}
\source{30}{889 名北美兽医：79\% 被要求提供自己认为无效的治疗}{Moses L, Malowney MJ, Wesley Boyd J (2018). Ethical conflict and moral distress in veterinary practice. Journal of Veterinary Internal Medicine 32(6):2115–2122. DOI 10.1111/jvim.15315（摘要经 Europe PMC 核实）}{Seventy-nine percent of participants report being asked to provide care that they consider futile.}{2026-09-18}{已核实}{反方论点二}{兽医终身暴露率，不是饲主比例}
\source{31}{477 名小动物兽医：99.2\% 职业生涯遇到过无效治疗，85.0\% 过去一年遇到过，42.4\% 认为常见（每年超过 6 次）}{Peterson NW, Boyd JW, Moses L (2022). Medical futility is commonly encountered in small animal clinical practice. JAVMA 260(12):1475–1481. DOI 10.2460/javma.22.01.0033（摘要经 Europe PMC 核实）}{85.0\% (402/476) reported encountering it within the past year}{2026-09-18}{已核实}{反方论点二}{同上}
\source{32}{同一调查作者评论文：76.2\% 兽医同意无效治疗有时有益于主人；继续无效治疗的理由包括满足主人“穷尽所有治疗手段”的要求}{Peterson NW, Boyd JW (2022-06-21). Veterinarians Often Provide Futile Care. Doing So Comes at a Cost. The Hastings Center，\url{https://www.thehastingscenter.org/veterinarians-often-provide-futile-care-doing-so-comes-at-a-cost/}}{76.2\% of respondents agreed or strongly agreed that providing futile care sometimes benefited the owners}{2026-09-18}{已核实}{反方论点二}{数字出自作者评论文，不在论文摘要中，引用时注明}
\source{33}{《2024 中国宠物体重状况调研报告》1,354 份：单犬家庭 69\%、单猫家庭 63\% 认为体态理想；被提示肥胖后 61\% 了解但无应对措施；犬猫超重肥胖 20\%–30\%；每周喂零食 3 次以上 31\%}{久生 Joyzone × 宠业家，2024-10-12；维宠宠物导航网转载，\url{https://www.petdhw.com/ganzhou-show-44771.html}}{61\%的宠物主表示了解，但并无应对措施；在单犬家庭中，69\%的宠主认为狗狗处于理想体态}{2026-09-18}{已核实}{反方日常战场}{转载页、宠粮品牌参与、在线问卷；61\% 的分母是全体饲主}
\source{34}{皇家全球调研（8 国、1.4 万宠主、1,750 兽医）：全球 40\% 成年犬猫超重或肥胖，中国 50\%；超过四分之一宠主对“健康体重”认知模糊}{皇家国际兽医研讨会新闻稿 2025-04-18，转载 \url{https://www.fabugov.com/html/jiankang/2025/0418/33037.html}}{全球40\%成年犬猫面临超重或肥胖问题；在中国，这一比例高达50\%}{2026-09-18}{已核实}{反方备用}{企业新闻稿转载，体况评估方法未公开，只作副证}
\source{35}{拟人化三因素理论：引发的施动者知识、效能动机、社交动机；缺乏与他人的社会连接时更倾向拟人化}{Epley N, Waytz A, Cacioppo JT (2007). On seeing human: A three-factor theory of anthropomorphism. Psychological Review 114(4):864–886. DOI 10.1037/0033-295X.114.4.864（摘要经 Europe PMC 核实）}{more likely to anthropomorphize ... when lacking a sense of social connection to other humans}{2026-09-18}{已核实}{反方日常战场}{}
\source{36}{英国 Vets at Home × Pet Database 宠物主调查：53\% 宠物死后感到内疚；确信时机正确者仍 47\% 内疚；25\% 对处理方式后悔；英国兽医诊所内 92\% 犬死亡、84\% 猫死亡为安乐死}{Companion Life 转载 2026-09-07，\url{https://www.companionlife.co.uk/new-research-reveals-pet-owners-unprepared-for-one-of-the-hardest-decisions-they-may-face/}}{More than half (53\%) of those surveyed said they felt guilt after their pet died}{2026-09-18}{已核实}{反方论点二}{样本量与抽样方法未公布，场上只说“英国一项调查”}
\source{37}{德牧“包包”恶性肿瘤安乐死后克隆，全程预计 30 万以上；猫“懒懒”患猫癣和肾积水后被家人扔在救助站附近}{青年志 Youthology《弃养、治疗、克隆，如何面对宠物的生老病死？》腾讯新闻 2023-10-13，\url{https://news.qq.com/rain/a/20231013A02XIW00}}{预计的全程花费在30万以上；扔在了离家很远的一家救助站附近}{2026-09-18}{已核实}{反方结辩画面、案例}{个案；同文的“流浪动物 4000 万”见第 43 条}
\source{38}{中国宠物医疗费用量级：狗跟腱断裂手术 5500 元、猫 CT 8000 元}{HUST 数说（华中科技大学新闻学院）《宠物看病比人还贵？》澎湃号 2025-11-28，\url{https://m.thepaper.cn/newsDetail_forward_32037301}}{一只狗狗的跟腱断裂手术费高达5500元，给猫猫做一次CT治疗要花费8000块}{2026-09-18}{已核实}{双方（账单真实）}{学生数据新闻，费用为个案汇总}
\source{39}{PDSA 2024：英国 19\% 的狗在典型工作日被独自留 5 小时以上，26\% 从不被单独留下}{PDSA Animal Wellbeing (PAW) Report 2024「Dogs」，\url{https://www.pdsa.org.uk/what-we-do/pdsa-animal-wellbeing-report/paw-report-2024/dogs}}{In 2024, 19\% of dogs are left alone for five or more hours}{2026-09-18}{已核实}{备用（未进稿件）}{}
\source{40}{照护负担量表 ZBI：负担是照护者对照护如何影响自己生活的主观评价}{Zarit SH, Reever KE, Bach-Peterson J (1980). Relatives of the impaired elderly: correlates of feelings of burden. The Gerontologist 20(6):649–655}{}{}{有把握}{反方论点二学理}{赛前确认卷期页码与量表定义原句}
\source{41}{照料行为系统与依恋系统是两套目标不同的系统}{Mikulincer M, Shaver PR (2005). Attachment Security, Compassion, and Altruism. Current Directions in Psychological Science 14(1):34–38}{}{}{有把握}{正方反驳库}{SagePub 403；赛前确认}
\source{42}{宠物作为“社会寄生者”破解人类育幼反应；拟人化选择塑造伴侣动物外形}{Archer J (1997). Evolution and Human Behavior 18(4):237–259；Serpell JA (2003). Society \& Animals 11(1):83–100}{}{}{有把握}{反方反驳库}{付费墙；赛前查校内数据库确认摘要}
\source{43}{中国流浪动物约 4000 万只}{青年志 2023-10-13 文中转述，出处未标；另有媒体转述《2021 年中国宠物行业白皮书》}{}{}{待核实}{未进稿件}{查中国小动物保护协会或农业农村部是否有官方估计；查不到就场上不用}
\source{44}{给狗穿衣“不到一小时可致中暑”这一具体时长}{见于 Mota-Rojas 2021 综述正文的二手转述}{}{}{待核实}{未进稿件}{MDPI 正文 403；打开正文后确认具体表述，之前只说“可致中暑甚至死亡”}
\source{45}{“单只犬年均消费 3006 元已超过全国居民人均看病买药支出”}{新浪财经 2026-09-14 报道的比较}{}{}{待核实}{未进稿件}{核对国家统计局 2025 年居民人均医疗保健支出后再用}
\source{46}{极端短头品种犬平均寿命 8.6 年，其他品种 12.7 年}{多篇科普转述}{}{}{待核实}{未进稿件}{查犬种寿命队列研究原文}
\end{sourcelist}
\textbf{统计}：已核实 39 条 / 有把握 3 条 /【待核实】4 条。
\textbf{给队员的硬规矩}：标【待核实】的四条，在有人查出原始出处之前，\textbf{一律不进任何一篇稿子、不在场上报数}。三条「有把握」只出现在学理与反驳库里，场上只说理论名与提出者，不报卷期页码。
\section{十二、备赛建议}
\begin{fields}
\begin{timeline}[时间表 （按距比赛 10 天倒排）]
\when{第 10–8 天}{全队通读本文档第三、四节，把“冲突时优先谁”和双方义务背到能用自己的话讲；每人写一遍“什么算为我 / 什么算为它”的三条。分工去查第十一节的四条【待核实】。}
\when{第 7–5 天}{稿件成型。正一稿已定，其余环节从口径表取词；二辩与四辩把第七节的被盘问预案逐条对答一遍。}
\when{第 4–3 天}{两场模辩，正反各打一次。重点检验下面三件事。}
\when{第 2–1 天}{按复盘调整，压时间，出声念稿。最后一遍跑 \texttt{check\_speeches.py}、\texttt{check\_voice.py}、\texttt{check\_consistency.py}、\texttt{check\_evidence.py}。}
\end{timeline}
\field{模辩重点检验的三件事}{
\begin{enumerate}[leftmargin=1.8em,itemsep=2pt,topsep=2pt]
\item 正方二辩与四辩对“负担高就考虑安乐死怎么解释”这一问的答案是否一字不差：开关在它的指标上。
\item 正方有没有在任何一处用 84\%、91\%、65\% 证“多数”，或说出“撑到撑不住才停”“被需要不是我的需要”。
\item 反方有没有在前 30 秒说“我方不否认你爱它”，有没有在任何一处说“养宠的人自私”“别骗自己了”。
\end{enumerate}
}
\begin{checklist}[每位队员赛前必须背下来的]
\item 我方的中立定义三句（爱、更、它/我）。
\item 中立判准“冲突时优先谁”，以及“我方要证明什么、对方要证明什么”。
\item 我方论点的标签与一句话主张；停止开关那一句（正方：在它的指标上；反方：跟着我的负担和不舍走）。
\item 对方论点的标签与我方的一句话反驳。
\item 我方全场必问的那一个问题（正方：什么才算爱的更是它；反方：同一份报告你引哪一题）。
\item 三个数字及其出处：本方口径表“数据”栏的前三行。
\end{checklist}
\end{fields}
\end{document}
